% 第 56 章　原子核与放射性（补充篇）
\chapter{原子核与放射性（补充篇）}

第~51~章里，卢瑟福用一束 \(\alpha\) 粒子轰击金箔，发现原子的几乎全部质量都缩在一个只有原子大小十万分之一的核里。此后五章我们一直绕着这个核打转：它提供正电荷，拴住电子云，决定了元素周期表和整个化学。但核本身被我们当成了一颗不会变化的带电小球。本章要钻进这颗小球里面看看。结果会出人意料：原子核不是一个静态的零件，而是一锅时刻可能重新洗牌的汤——有的核会自发地换一种身份，同时吐出携带能量的碎片。这就是放射性。它是炼金术士梦寐以求的“元素互变”，只不过发动它的不是巫师，而是概率。

\step{反常识开场}

先看一个你可能从没意识到的“事实”：此刻，你自己的身体就是一个放射源。你体内大约有一百四十克钾，其中极小的一撮是钾的一种不稳定的孪生兄弟——钾-40；它们之中每秒约发生四千次衰变，每一次都向你的组织里射出一个高速电子或一个高能光子。再加上你骨骼和脂肪里的碳-14，每秒再添三千次左右。合计下来，你的身体内部每秒发生七八千次“微型核爆炸”，一年超过两千万亿次。你没有发光，没有发热到异常，也毫无感觉。

再听一个更老的怪事。一八九六年，贝克勒尔把铀盐放在包好的照相底片上，原本想研究荧光，结果天阴没做成实验，铀盐在抽屉里躺了几天——底片却自己曝光了，位置正好是铀盐躺过的地方。没有光照，没有通电，没有任何人碰它，一块石头在黑暗里持续不断地向外发射某种能穿透黑纸的“东西”。后来居里夫人发现，这种发射跟温度无关，跟化学处理无关，跟放多久无关：你把铀盐磨碎、溶解、冻住、烧红，它发射的劲头一丝不变。一种完全不听外界使唤、永远关不掉的“机关”，藏在一块石头里。

最后一个悬念藏在你家的天花板上。许多烟雾报警器里装着一粒纽扣大小的放射性金属镅-241，它就挂在你头顶，日夜不停地放出 \(\alpha\) 粒子。放射性不是核电站里的稀罕物，它在你的身体里、你的房子里、你的土壤里。问题是：原子核为什么要“衰变”？它什么时候衰变？放出来的东西去了哪里？以及——你每天被自己轰击七八千次，为什么还活着？

\step{先猜一猜}

照旧，先把你的直觉写在纸上，再往下读。

第一猜：原子核里全是带正电的质子（后面还会讲到中子，但它不带电）。同性电荷互相排斥，而它们被挤在半径只有一飞米（\(10^{-15}\)~m）量级的空间里。你猜拴住它们的是（a）万有引力，（b）电磁吸引力的某种残余，（c）一种你还没见过的新的相互作用？可以先算一笔账再猜：第~19~章的库仑定律说，两个质子相距一飞米时，互相推开的电势能大约是一百四十万电子伏特；而万有引力在同距离下提供不了它的万亿分之一。

第二猜：一块一克的镭，天然地、不停地放出射线。你猜它每秒大约衰变多少个原子？是几百个、几百万个，还是更多？再猜：如果把这块镭埋起来，等一千年，它放出的射线会减弱大约多少？

第三猜，也是最要紧的一猜：假设你手里有一百个同一种放射性原子，已知这种原子的“半衰期”是一天，也就是一天之后平均剩五十个。那么两天之后剩二十五个，三天之后剩十二三个……请问第几天它们会“死光”？更根本的问题是：盯住其中任意一个原子，你能预测它明天衰不衰变吗？它已经等了一整天没有衰变，是不是比刚造出来时“更累”“更接近死期”了？

\step{动手实验}

\begin{experiment}[用一百枚硬币模拟衰变]
\textbf{材料}：一把硬币（或纽扣、豆子、骰子，总数五十到两百之间），一张纸，一支笔。

\textbf{第一步}：数一数你的硬币总数 \(N_0\)，记在表格第一行。把硬币撒在平面上，把所有正面朝上的挑出来拿走——它们“衰变”了。数一数还剩多少，记下“第 1 轮之后剩余 \(N\)”。

\textbf{第二步}：把剩下的硬币收拢，再撒一次，再挑走正面朝上的。如此重复八到十轮，每轮记录剩余数。

\textbf{第三步}：在坐标纸上画出剩余数随轮数变化的点，再在旁边画一条“每轮精确减半”的曲线（从 \(N_0\) 出发，依次取 \(N_0/2\)、\(N_0/4\)、\(N_0/8\)……）对比。

\textbf{记录与思考}：你的点大致贴着减半曲线走，但每一轮都有偏差。请回答三个问题。其一：你能提前说出某一枚特定的硬币下一轮会不会被拿走吗？其二：为什么一百枚硬币的整体行为，比十枚硬币的整体行为更“守规矩”？其三：如果某一枚硬币连续五轮都“幸存”了，第六轮它被挑走的概率变大了吗？请把你的答案留到“回头解释开场现象”一节，和真正的原子核对答案。
\end{experiment}

这个实验是理解本章全部统计规律的钥匙，请务必真的做一遍，哪怕只用三十枚硬币。它和真实放射性的对应关系是：每一轮代表一个固定的时间间隔，每枚硬币代表一个原子核，“正面朝上就出局”代表“在这段时间里衰变了”。实验里那个让你不舒服的事实——单枚不可预测、整体却极其规律——正是放射性最深刻的性质，不是我们的仪器不够好造成的，而是原子核的衰变在原则上就没有倒计时。

\step{旧模型遇到麻烦}

现在用第~19~章的电磁学和第~52~章的原子图像，认真解释原子核，看它在哪里翻车。

第一场翻车：原子核根本不该存在。以氦核为例：两个带正电的质子挤在约一飞米的距离内，库仑排斥对应的电势能约 \(1.4\times10^{6}\)~eV。给这个能量一个直观标尺：要让热运动的动能达到这个量级，温度得超过 \(10^{10}\)~K——太阳中心的温度也才 \(1.5\times10^{7}\)~K，差着近千倍。万有引力呢？同距离下两个质子之间的引力势能只有库仑排斥能的约 \(10^{-36}\) 分之一，可以彻底无视。按旧模型，原子核应该像一小把同极相对的磁铁，在飞米尺度上被一股相当于“万亿度高温”的排斥力瞬间炸开。可是铀核在岩石里安安稳稳坐了几十亿年。旧模型预言了灾难，大自然却毫发无损——这不是误差，是缺了整整一种相互作用。

第二场翻车：质量和电荷对不上账。称一称原子核：氦核的质量约是氢核的四倍，电荷却只有两倍；碳核质量是氢核的十二倍，电荷只有六倍。如果核里只有质子，质量与电荷应该永远成正比，可它们偏偏不是。历史上有人打过补丁：假设核里除了质子还有电子，电子中和掉多余的质子电荷，质量却不增加多少。这个“核内电子”补丁越打越漏：按第~44~章的不确定性关系估算，把电子关在飞米尺度的核里，它的动量不确定度大到对应几十 MeV 的能量，核里根本关不住这么轻的东西；而且对核自旋的测量也与这个补丁矛盾。补丁的葬礼在一九三二年举行：查德威克从铍的辐射里找到了一种不带电、质量与质子几乎相同的新粒子——中子。账立刻平了：氦核是两个质子加两个中子，质量四倍、电荷两倍，分毫不差。

第三场翻车：\(\beta\) 衰变似乎在撕毁能量守恒。某些核会放出一个电子（这叫 \(\beta\) 射线），核的电荷数因此加一。麻烦在于：如果每次衰变都是同一个核变成同一个新核，那么放出的电子应该带走一个确定的能量，就像第~51~章里 \(\alpha\) 粒子的能量那样整齐。实验却给出一张连续的能谱：电子的能量从接近零到某个最大值都有，每次都不一样，而且电子的反冲方向与核的反冲方向也对不上账。要么能量守恒和动量守恒在原子核里失效，要么——有一个看不见的第三者把剩下的能量和动量带走了。泡利在一九三〇年押了后者：他假设存在一种不带电、几乎不与任何东西相互作用的粒子，偷偷分走了账目的余额。这个“做账平账”的粒子后来被命名为中微子，二十六年后才第一次被直接探测到。这一次，守恒律没有输，输的是“账本里只有看得见的粒子”这个旧假设。

\step{发明新概念}

三场翻车指向同一张设计图。像历史上的物理学家一样，我们把新东西一件件发明出来。

第一样：\textbf{核子}。质子和中子统称核子。它们在核里的地位几乎平等：质量相差不到千分之一点四，都不怕彼此的电排斥——中子干脆不带电。一个原子核由它的质子数 \(Z\)（决定它是哪种元素、核外配多少电子、化学脾气如何）和中子数 \(N\) 共同标记，核子总数 \(A=Z+N\) 叫质量数。质子数相同、中子数不同的核互为\textbf{同位素}：化学家分不出它们（核外电子一模一样），核物理学家却分得清清楚楚——氢、氘、氚就是这样一家三兄弟，前两个安稳，第三个有放射性。

第二样：\textbf{强相互作用}。核不散架，是因为核子之间存在一种全新的吸引：它极强，在近一飞米的距离上轻松压过质子间的库仑排斥；它又极短程，超过两三飞米就几乎归零——所以它从不出现在化学和日常生活里，你在针筒里、在桌面上永远感觉不到它。它还有一个朴素的性格：与电荷无关，质子—质子、中子—中子、质子—中子之间的吸引力差不多一样大；并且每个核子只与紧邻的少数几个核子“拉手”（这叫饱和性），不像电荷那样隔空对所有带电粒子一视同仁。这三条性格合起来，决定了原子核的全部宏观面貌。

第三样：\textbf{结合能与质量亏损}。把一个核拆成单个核子需要花能量，这份能量叫结合能 \(B\)。它从哪里记账？第~40~章的质能关系给出答案：拆开时吸收的能量会变成质量。所以一个结合着的核，它的静止质量 \(M\) 严格小于各核子自由质量之和，差额叫质量亏损：
\[
B=\bigl(Zm_{\mathrm p}+Nm_{\mathrm n}-M\bigr)c^{2}.
\]
用日常语言说：把核子粘在一起后，整碗汤比所有原料分开称还要轻，轻掉的那部分质量就是胶水里的能量。这不是“质量变成了能量”的玄学，而是能量本身有惯性、上秤会显示出质量的直接证据。这里说的质量始终指静止质量（不变质量），是核本身的固有属性。

第四样：\textbf{三种衰变}。不稳定的核通过三条途径换个更稳的身份。\textbf{\(\alpha\) 衰变}：核抛出一个氦-4 核（两个质子加两个中子），于是 \(A\) 减四、\(Z\) 减二，元素变了一种。\textbf{\(\beta\) 衰变}：核内一个中子转变成一个质子（或反过来），同时放出一个电子（或正电子）和一个中微子，\(A\) 不变、\(Z\) 增减一，元素也变了——掌管这种“核子换身份”操作的，是继引力、电磁、强相互作用之后的第四种基本相互作用，叫弱相互作用。\textbf{\(\gamma\) 衰变}：核从高能态跳回低能态，放出一个能量极高的光子，核的成分一点都不变，只是“消了消气”。三种衰变的穿透本领天差地别：\(\alpha\) 粒子连一张纸和皮肤表层都穿不过，\(\beta\) 电子能被几毫米铝板拦住，\(\gamma\) 光子则要厚厚的铅板或混凝土才挡得住——这一条后面谈安全时要反复用到。

第五样，本章的心脏：\textbf{衰变常数与半衰期}。实验事实是：每一种不稳定核的每一个个体，在单位时间内都有一个固定的衰变概率，记作 \(\lambda\)，而且这个概率与它活了多久、同伴还剩多少、外界温度压强如何，全都无关。请注意措辞：这是关于\emph{单个}原子核的概率陈述，不是“到了某时刻集体赴死”的倒计时。由此推出一个整体规律：一大群这样的核，平均每隔一段固定时间 \(T_{1/2}\)（半衰期）就衰变掉一半。半衰期是统计规律，是大数定律在原子核上的显影；单个原子核没有年龄、没有履历、没有“死期将近”——这一点，你的硬币实验已经提前告诉你了。

\step{一句话模型}

\begin{oneline}
原子核是质子和中子靠短程的强相互作用捆成的团：捆得紧不紧，由比结合能（每个核子分摊的结合能）决定，铁附近捆得最紧；捆得不够紧的核会以固定的概率随机“换身份”——放出 \(\alpha\)、\(\beta\) 或 \(\gamma\) 射线——于是一大群核的数量按半衰期指数衰减，而单个核的衰变时刻永远不可预测。
\end{oneline}

这个模型像什么？像一间挤满人的候车厅：每个人都揣着一枚自己的骰子，每隔固定时间掷一次，掷出某个点数就起身离场（衰变），离场的概率人人相同、轮轮相同。厅里人越多，你预测“这一轮大约走多少人”就越准，可你永远说不出张三第几轮走。

它不像什么？有三处必须立刻标明。第一，候车厅里的人会累、会饿、会熬不住，离场概率随等待时间上升；原子核不会——一个已经等了十亿年的铀核和一个刚诞生的铀核，下一秒衰变的概率一模一样，这叫“无记忆性”，直觉最容易在这里栽跟头。第二，人离场有原因（车来了），核衰变没有更深一层的“触发器”可查：不是仪器不够精细，而是量子理论给出的就是一个概率事件，与第~44~章隧穿、第~45~章概率同源。第三，骰子的点数是固定规则，而“哪种核、以多大的 \(\lambda\) 衰变”没有万能旋钮——它由核内强作用、弱作用和库仑排斥三方的精细平衡决定，改不了、调不动，这就是贝克勒尔那块石头“永远关不掉”的原因。

\step{数学翻译器}

先算结合能这笔真账，用真实数据。氦-4 核含两个质子、两个中子。质子的静止质量 \(m_{\mathrm p}=1.007276\)~u，中子 \(m_{\mathrm n}=1.008665\)~u，氦-4 核的实测质量 \(M=4.002603\)~u（u 是原子质量单位，\(1~\mathrm{u}\,c^{2}=931.5\)~MeV）。质量亏损：
\[
\Delta m=2\times1.007276+2\times1.008665-4.002603=0.030377~\mathrm{u},
\]
\[
B=\Delta m\,c^{2}=0.030377\times931.5\approx28.3~\mathrm{MeV},
\qquad \frac{B}{A}\approx\frac{28.3}{4}\approx7.1~\mathrm{MeV}.
\]
四个问题走一遍。描述什么：把氦核拆散需要的能量，以及每个核子平均分摊多少。为什么这样写：质能关系说能量与质量是同一笔账的两种记法，所以“轻了多少”乘上 \(c^2\) 就是“粘得多牢”。何时成立：只要各质量指静止质量，它精确成立，是实验事实而非模型。何时失效：它只给总账，不告诉你核内部结构，也预测不了哪个核会衰变。特殊检查：\(Z\) 或 \(N\) 取零，\(B=0\)——单个质子谈不上结合，对；质量亏损的符号若算成负的（核比原料还重），这个核根本不会被大自然组装出来，对——自然界只造 \(B>0\) 的核。

把每个核子的比结合能 \(B/A\) 对所有元素画出来，得到本章最重要的一张图：从氢出发一路上坡，氦-4 处有一个尖峰，在铁（\(A\approx56\)）附近到达顶点（约 \(8.8\)~MeV），随后缓缓下坡，到铀（约 \(7.6\)~MeV）。这条曲线是核世界的“地势图”：任何让核朝铁靠拢的变动——轻核聚合（聚变）或重核分裂（裂变）——都是沿比结合能“上坡”（也就是总能量“下坡”），都会放能。恒星燃烧和核电站发电，用的是同一张图的两个斜坡。

\begin{figure}[htbp]
\centering
\begin{tikzpicture}[>=Stealth,scale=1.0]
  \draw[->] (0,0) -- (10.2,0) node[below left=1pt]{核子数 \(A\)};
  \draw[->] (0,0) -- (0,5.6) node[left=1pt]{\(B/A\)（MeV）};
  \foreach \y/\lab in {1/2,2/4,3/6,4/8}{
    \draw (-0.08,\y*0.6+0.6) -- (0.08,\y*0.6+0.6) node[left=4pt]{\scriptsize \lab};
  }
  % 比结合能曲线：起点氢(0)，快速上升，Fe 峰，缓慢下降
  \draw[thick,blue!60!black] plot[smooth] coordinates {
    (0.15,0.6) (0.5,3.9) (1.0,4.4) (2.0,4.9) (3.5,5.25) (4.5,5.3)
    (5.5,5.2) (7.0,4.95) (8.5,4.75) (9.7,4.6)};
  \fill (4.5,5.3) circle (1.5pt) node[above=2pt]{\scriptsize 铁 \(A\approx56\)};
  \fill (0.5,3.9) circle (1.5pt) node[above right=0pt]{\scriptsize 氦-4};
  \fill (9.7,4.6) circle (1.5pt) node[above=1pt]{\scriptsize 铀};
  \draw[->,red!60!black,thick] (1.2,5.5) -- (3.4,5.5) node[midway,above]{\scriptsize 聚变：轻核上坡放能};
  \draw[->,red!60!black,thick] (9.0,5.9) -- (6.2,5.9) node[midway,above]{\scriptsize 裂变：重核上坡放能};
\end{tikzpicture}
\caption{比结合能曲线：铁在最“高”处捆得最紧，轻核聚变与重核裂变都在向铁靠拢，同时放出能量。}
\end{figure}

再翻译衰变的统计规律。主线层先给结论：若某个时刻有 \(N_0\) 个相同的放射性核，经过时间 \(t\) 后平均剩余
\[
N(t)=N_{0}\times2^{-t/T_{1/2}} .
\]
特殊检查立刻做：\(t=0\) 给出 \(N_0\)，对；\(t=T_{1/2}\) 给出 \(N_0/2\)，正是半衰期的定义；\(t=2T_{1/2}\) 给出 \(N_0/4\)；\(t\) 趋于无穷大时 \(N\) 趋于零却永远不等于零——所以在第三猜里问“第几天死光”，答案是：平均意义上永远死不光，只是少到看不见；而“一百个原子的最后一天”根本没有确定的日子。

\begin{mathbridge}
为什么偏偏是指数函数？把“每个核在单位时间内衰变的概率都是 \(\lambda\)”翻译成一句话：在极短时间 \(dt\) 内，平均有 \(\lambda N\,dt\) 个核衰变，于是剩余数的变化满足
\[
\frac{dN}{dt}=-\lambda N .
\]
右边的 \(N\) 是关键：衰变速率正比于还剩多少——核越少，衰变得越慢。这个方程的解是
\[
N(t)=N_{0}e^{-\lambda t}.
\]
令 \(N=N_0/2\) 解出半衰期：\(T_{1/2}=\ln2/\lambda\approx0.693/\lambda\)。两条式子完全等价，\(2^{-t/T_{1/2}}\) 就是 \(e^{-\lambda t}\) 换了个底。再验证无记忆性：一个核已经活了时间 \(s\) 还没衰变，它再活 \(t\) 的概率是 \(e^{-\lambda(s+t)}/e^{-\lambda s}=e^{-\lambda t}\)——与 \(s\) 无关。等多久都白等，概率清零重来。
\end{mathbridge}

每秒的衰变次数叫活度 \(A=\lambda N\)，单位是贝克勒尔（Bq），一 Bq 就是每秒一次。现在兑现开场的第一笔估算，全程真实数据。七十公斤的成年人含钾约 \(140\)~g，其中钾-40 的天然丰度约为 \(0.0117\%\)，即约 \(0.0164\)~g，合 \(N=0.0164/40\times6.0\times10^{23}\approx2.5\times10^{20}\) 个原子。钾-40 的半衰期 \(T_{1/2}=1.25\times10^{9}\) 年，折合约 \(3.9\times10^{16}\)~s，于是
\[
A=\lambda N=\frac{\ln2}{T_{1/2}}N\approx\frac{0.693\times2.5\times10^{20}}{3.9\times10^{16}}\approx4.4\times10^{3}~\mathrm{Bq}.
\]
每秒四千四百次，正是开场说的数字。再用同一公式给居里夫人算一笔：一克镭-226 有 \(N=6.0\times10^{23}/226\approx2.7\times10^{21}\) 个原子，半衰期约 \(1600\) 年，活度约 \(3.7\times10^{10}\)~Bq——历史上放射性活度的旧单位“居里”就是这么定义的。顺便回看第二猜：每秒三百七十亿次，比“几百万”大了四个数量级；一千年（约零点六个半衰期）后，射线强度大约还剩三分之二。放射性世界的数字，永远比直觉给出的夸张。

活度只说“每秒衰变多少次”，不直接等于伤害。完整的安全账要分三个量：活度（Bq）数衰变；吸收剂量（Gy，戈瑞）称每千克组织吸收了多少焦耳能量；当量剂量（Sv，希沃特）再按射线种类加权——同样一焦耳能量，\(\alpha\) 造成的生物损伤约为 \(\gamma\) 的二十倍，所以乘上不同的权重因子。天然本底辐射（宇宙线、土壤、氡气、你体内的钾）平均每人每年约 \(2.4\)~mSv；一次胸片约 \(0.1\)~mSv；一次胸部 CT 约 \(7\)~mSv；一年累计 \(100\)~mSv 以下，流行病学上测不出可分辨的额外风险。这三层账各司其职，混用它们是绝大多数辐射恐慌的源头。

最后是能源账，还是真实数据。一个铀-235 核裂变平均放出约 \(200\)~MeV，即 \(3.2\times10^{-11}\)~J。一克铀-235 含 \(2.6\times10^{21}\) 个核，全部裂变放出约 \(8.2\times10^{10}\)~J——相当于约 \(2.8\) 吨标准煤，或一户家庭十几年的用电量。一克换将近三吨，这就是质能关系在比结合能曲线上兑现的汇率。太阳那头更夸张：它每秒辐射 \(3.8\times10^{26}\)~J，按 \(E=mc^2\) 折算，相当于每秒把约 \(420\) 万吨静止质量转成辐射撒向太空；它已经这样烧了四十六亿年，靠的正是曲线上从氢到氦那一段上坡。

\step{模型的边界}

强相互作用的图像说到哪算数？本章把它当作“飞米尺度上的强吸引力”是有效的描述，但它不是一根能画出弹性系数的弹簧：它随距离剧烈变化，在两飞米外迅速归零，在更近处还会反转为排斥（否则核子会塌成一个点）。再往下挖一层，质子和中子本身是由夸克组成的复合粒子，强相互作用是夸克之间胶子相互作用的“残余”——就像分子间的范德华力是电磁相互作用的残余。这是第~59~章“相互作用的地图”里要展开的题材，本章只需记住：我们的“强力”是有效模型，在飞米尺度上算账好用，出了这个尺度就别外推。

比结合能曲线是一张地势图，不是万有公式。它画的是平滑趋势，真实的核在这条趋势线上叠加了壳层结构的起伏：质子和中子在核里也按类似电子壳层的方式排布，某些“幻数”组合（如 \(2,8,20,28,50,82,126\)）格外稳定。曲线预测不了哪个具体的核会以哪种方式衰变——那是壳层细节、库仑排斥与弱相互作用三方谈判的结果。

半衰期这条统计规律的适用边界更要刻在脑子里。它要求原子数巨大：十亿个原子谈“一半衰变”意义精确；十个原子谈“半衰期”就只剩平均值的虚名，你的硬币实验里那一上一下的涨落就是证据。单个原子核没有半衰期可言，只有概率。反过来，指数律在极端情形下也会失真：当剩余原子少到个位数，衰变时刻完全由涨落主宰；而在比 \(10^{-23}\) 秒还短的瞬间，量子理论预言衰减尚未进入指数区。这些修正对日常和实验室都无关紧要，但它们提醒你：指数律是模型，不是铁的教条。

\(\beta\) 衰变那张连续能谱留给我们的方法论教训，值得单独画一条边界：当年“能量守恒在核里失效”的诊断错得彻底，真相是账本上漏了一个几乎看不见的粒子。守恒律的正确用法不是“看起来不平就宣布失效”，而是“不平就去找漏记的项”。中微子是这条方法论最辉煌的战利品——但也请注意，发现它的不是会计技巧，而是二十六年后实打实的探测：假设必须最终兑成测量。

辐射安全里有两条广泛流传的误解，本章有责任澄清。第一，“被照过就变成放射源”。错。让普通物质带上放射性需要用中子轰击使其核发生改变（这叫活化）；X 光机、\(\gamma\) 源照射过的病人和食物，其原子核成分没有改变，不会因此自己发射射线——辐照食品离开辐照室那一刻，放射性就随射线一起走了。第二，“辐射剂量越低越好、任何剂量都有可测的伤害”。低于每年约 \(100\)~mSv 的剂量，额外风险小到今天无法从统计噪声里分辨；“线性无阈”外推（假设风险随剂量一路线性降到零）是一种审慎的管理模型，不是已被证实的实验事实——把它当事实，正是“香蕉恐慌”与“体检 CT 恐慌”共同的来源。另外，天然辐射里最大的单一来源不是核电站也不是医疗，而是地基渗进室内积存的氡气：通风，是普通人对辐射能做的最实际的一件事。

\begin{challenge}
为什么不同核的半衰期跨度如此恐怖——铀-238 是 \(45\) 亿年，钋-212 是 \(0.3\) 微秒，相差 \(10^{24}\) 倍，而它们放出的 \(\alpha\) 粒子能量只差一倍多一点（约 \(4.3\)~MeV 对 \(8.8\)~MeV）？一九二八年伽莫夫给出的解释是第~44~章隧穿效应的真人秀：\(\alpha\) 粒子在核内被强作用的势阱扣住，外面顶着库仑排斥的势垒，经典物理里它永远翻不出去，量子力学却允许它以概率 \(P\approx e^{-2G}\) 穿垒而出，其中 \(G\) 对势垒的高度和宽度极其敏感。\(\alpha\) 能量每高一点，势垒就显得矮一点、薄一点，穿透概率按指数放大——能量只差一倍，概率能差二十多个数量级。一条指数关系，把“亿年”和“微秒”缝在同一张图上。这也解释了为什么 \(\alpha\) 衰变是量子隧穿在人类眼皮底下最早的一次亮相。
\end{challenge}

\step{回头解释开场现象}

现在把开场的三件事一件件收回来。

你体内的每秒七八千次衰变，来自钾-40 和碳-14 这两种天然放射性同位素。它们为什么赶不走？因为同位素的化学性质由核外电子决定，你的肾脏和细胞膜分不出钾-39 和钾-40——对化学而言它们是同一种元素，吃进去的钾里永远按比例掺着那一小撮放射性兄弟。这不是污染，是地球化学的出厂配置：四十亿年前形成地球的原料里就有钾-40，它比地球活得还从容（半衰期十二点五亿年），生命从诞生第一天起就泡在这种本底里。至于你为什么扛得住：每次衰变沉积的能量微乎其微，且分散在全身几十万亿个细胞上；细胞自带一套修复损伤的分子机器，对每年约一毫希沃特量级的内照射应付裕如——这本身就是演化的结果，因为从来没有一个“无辐射”的环境可供生命选择。

贝克勒尔的铀盐为什么永远关不掉？因为衰变概率 \(\lambda\) 写在核内强作用、库仑排斥与弱相互作用的平衡里，而温度、压强、化学状态全是核外电子层面的事，能量尺度差着几十万倍——你用火烧、用酸泡，折腾的都是电子云，核内那锅汤连感觉都没有。一八九六年那张自己曝光的底片，是铀-238 及其衰变子孙放出的 \(\beta\) 与 \(\gamma\) 穿透黑纸打出来的；而铀的半衰期长达四十五亿年，所以那块石头今天还在以几乎相同的劲头发射着。

烟雾报警器则是把放射性驯化成工具的经典案例。镅-241 放出的 \(\alpha\) 粒子把探测器小腔室里的空气分子电离，维持一股微弱而恒定的电流；烟雾颗粒飘进腔室后会吸附离子，电流一下降，警报就响。为什么敢把它挂在卧室？因为 \(\alpha\) 粒子连几厘米空气和一层塑料外壳都穿不过，只要不拆开吞下去，剂量可以忽略。同一种射线，关在体外是纸都穿不过的“软柿子”，被吸入或吞进体内却会贴着细胞释放全部能量——氡气的危害正来自这条不对称：危险与否，取决于放射源在围墙的哪一边。

裂变与聚变则是比结合能曲线上的两次“上坡冲刺”。铀-235 吸收一个中子后裂成两块更靠近铁的碎片，比结合能上坡，差额变成碎片的动能；每个裂变还顺带放出两三个中子，去敲开旁边的铀核——链式反应。核电站的全部技术，慢化剂、控制棒、厚厚的安全壳，都是围绕“让中子流不多不少”这一件事展开的精细管理。聚变反过来：两个氢的同位素聚成氦，比结合能从零附近冲上七点一的尖峰，放能更多；但两个带正电的核要先顶着库仑排斥贴到一飞米内——太阳中心一千五百万度都不够“烧”过去，靠的是第~44~章的隧穿：质子不必翻过势垒，只需穿过它，聚变速率才刚好够太阳烧上几十亿年而不炸成烟花。人类在地球上复刻这团火（托卡马克、激光聚变），难就难在要把一亿度的“汤”悬空端住。而太阳、核电站、烟雾报警器和你体内的钾-40，在比结合能曲线这张图上是同一件事的四个注脚。

\step{脑洞实验与挑战题}

\begin{thought}[如果强相互作用的力程变成十倍]
做一个极端尺度的思想实验：假设强相互作用的有效力程从约一飞米变成十飞米，其余不变。会发生什么？首先，库仑排斥是长程的、随质子数一路累积，而强力原本是“只拉邻居”的短程手；力程放大十倍后，每个核子能拉住的邻居从几个变成几百个，吸引大幅反超排斥，重核会变得远比今天稳定——比结合能曲线的铁峰向更重的方向大幅移动，铀可能不再是天然存在的最重元素，地壳里会出现今天根本不存在的超重稳定元素。反过来，如果强力强度减半或力程缩到十分之一飞米，氦-4 都可能聚不住，恒星连第一把火都点不着，宇宙里几乎只剩氢。碳基生命要求的参数窗口窄得惊人：强一点、弱一点，周期表都要重写。这个实验没有旋钮可供验证（自然界不给我们这只旋钮），它的价值在于让你看清：我们世界的元素清单，是几种相互作用的强度与力程之间精细平衡的结果，而不是理所当然的默认值。
\end{thought}

\begin{puzzles}
\item 一块含一百万个某种放射性原子的样品，半衰期为一小时。大约多少小时后，平均剩余的原子不到一个？“不到一个”是什么意思？此时再说“半衰期”还有意义吗？（思路提示：每过一小时数量减半，计算 \(2^n\) 何时超过一百万；\(2^{20}\approx10^{6}\)。剩余数接近个位数时，衰变完全由涨落主宰，“半衰期”作为统计规律的前提——大数——已经垮掉，只剩单原子概率还有意义。）
\item 深空探测器常用钚-238（半衰期约八十八年）做核电池热源，而不用铀-235（半衰期约七亿年）或钋-210（半衰期约一百三十八天）。从“功率正比于活度 \(\lambda N\)”出发，说明为什么半衰期太长、太短都不合用。（思路提示：功率 \(=\lambda N\times\) 每次衰变能量。铀-235 的 \(\lambda\) 太小，同样质量活度低得可怜，发热功率微乎其微；钋-210 的 \(\lambda\) 很大，功率起初强劲，但几个月就烧掉一大半，撑不了十年的深空航行；八十八年恰好卡在“功率够用、又能撑几十年”的窗口里。）
\item 碳-14 测年法的原理是：活着的生物不断与大气交换碳，体内碳-14 比例与大气持平；死后交换停止，碳-14 以五千七百三十年的半衰期衰减。请问：为什么碳-14 可以测一件三千年前的木器，却测不了六千五百万年前的恐龙化石？（思路提示：三千年约为零点五个半衰期，剩约七成，信号充足；六千五百万年超过一万个半衰期，剩余比例是 \(2^{-10000}\)，连一个碳-14 原子都不会剩下——方法的前提“样品里还有可测的碳-14”彻底失效，必须换用半衰期更长的铀铅等时钟。）
\item \(\alpha\) 衰变放出的 \(\alpha\) 粒子能量几乎总是单一的某个值，\(\beta\) 衰变放出的电子能量却是一张连续谱。从“衰变产物有几方”出发解释这个差别，并说明当年它为何让物理学家怀疑能量守恒。（思路提示：\(\alpha\) 衰变是两体问题——母核变成子核加 \(\alpha\) 粒子，能量与动量守恒两条方程把两方的能量完全定死；\(\beta\) 衰变实际是三体——子核、电子，再加上当年没人知道的中微子，释放的总能量可以在三方之间任意分配，电子分到多少都行。漏看了第三者，账自然永远对不平。）
\end{puzzles}

本章的放射性故事到这里收线，但有三根线头有意留给了后面。其一，\(\beta\) 衰变背后那位“弱相互作用”，连同引力、电磁、强相互作用，将在第~59~章被放进同一张相互作用的地图，那里会问：这四种力为什么长得如此不一样，又为何可能写着同一种语言？其二，本章反复使用的能量守恒、动量守恒，在第~58~章会被追出更深的来历——它们不是实验的偶然总结，而是“规律不随时间、地点改变”这种对称性的直接产物。其三，衰变的随机性与第~45~章的量子概率一脉相承：单个原子核何时衰变，量子理论只给概率不给日程——这是理论的结构，不是我们测量的无能。原子核这锅汤里的每一场随机洗牌，洗的都是同一条规则的牌。
